% Trabalho Prático 01 --- documento consolidado de todos os domínios do PMBOK 8.
% O conteúdo de cada domínio NÃO é copiado aqui: os .tex de cada pasta são incluídos
% com \input via docmute (que ignora o preâmbulo de cada arquivo). Para mudar um
% domínio, edite o .tex da pasta dele e recompile este arquivo.
\documentclass[11pt,a4paper]{report}
\usepackage[utf8]{inputenc}
\usepackage[T1]{fontenc}
\usepackage[brazil]{babel}
\usepackage[margin=2.5cm]{geometry}
\usepackage{mathptmx}
\usepackage{amsmath}
\usepackage{amssymb}
\usepackage{enumitem}
\usepackage{booktabs}
\usepackage{tabularx}
\usepackage{xltabular}
\usepackage{array}
\usepackage{forest}
\usepackage{tikz}
\usetikzlibrary{arrows.meta,positioning}
\usepackage{docmute}
\usepackage{titlesec}
\usepackage{fancyhdr}
\usepackage{hyperref}
\hypersetup{colorlinks=false, pdfborder={0 0 0}, pdftitle={TrackFleet --- Documento Consolidado do Projeto}}

\pagestyle{fancy}
\fancyhf{}
\fancyhead[L]{\small TrackFleet --- Documento Consolidado do Projeto}
\fancyhead[R]{\small\thepage}
\renewcommand{\headrulewidth}{0.4pt}
\fancypagestyle{plain}{\fancyhf{}\fancyfoot[C]{\small\thepage}\renewcommand{\headrulewidth}{0pt}}

\titleformat{\chapter}[display]{\normalfont\bfseries}{\large\MakeUppercase{\chaptertitlename}\ \thechapter}{4pt}{\LARGE}
\titlespacing*{\chapter}{0pt}{-10pt}{18pt}

\setcounter{secnumdepth}{3}
\setcounter{tocdepth}{1}

\newcolumntype{L}{>{\raggedright\arraybackslash}X}
\newcolumntype{C}{>{\centering\arraybackslash}X}
\newcolumntype{R}{>{\raggedleft\arraybackslash}X}
\setlist{topsep=3pt}

% Comandos usados pelos documentos de cada domínio, redefinidos para o consolidado:
% o cabeçalho de documento some (o capítulo já identifica o domínio), cada parte de
% um domínio com vários documentos vira uma seção, e os separadores viram espaço.
\let\secaoparte\section
\newcommand{\cabecalho}[3]{}
\newcommand{\parte}[1]{\secaoparte{#1}}
\newcommand{\separador}{\par\medskip}
\newcommand{\dodbox}{$\square$}
% Inclui um documento que é uma das partes de um domínio: as seções dele descem um nível.
\newcommand{\subdocumento}[1]{\begingroup\let\section\subsection\let\subsection\subsubsection\input{#1}\endgroup}
% Capítulo sem número que ainda aparece no sumário.
\newcommand{\capitulosemnumero}[1]{\chapter*{#1}\addcontentsline{toc}{chapter}{#1}\markboth{#1}{#1}}
\newcommand{\secaosemnumero}[1]{\section*{#1}\addcontentsline{toc}{section}{#1}}

\begin{document}

%======================================================================
% Capa
%======================================================================
\begin{titlepage}
\centering
{\large SETREM --- Sociedade Educacional Três de Maio}\\[2pt]
{\large Engenharia da Computação $\cdot$ 4\textsuperscript{o} semestre}\\[2pt]
{\large Disciplina de Gestão de Projetos $\cdot$ Prof. Eduardo Mildner}

\vspace{5cm}

{\LARGE\bfseries TrackFleet}\\[6pt]
{\Large\bfseries Gestão de Rastreador de Automóveis}\\[24pt]
{\large DOCUMENTO CONSOLIDADO DO PROJETO}\\[4pt]
{\large Trabalho Prático 01}\\[16pt]
Framework PMBOK\textregistered\ 8\textsuperscript{a} Edição $\cdot$ os sete domínios de desempenho

\vfill

{\large Enzo Allebrand e Kerlon Ribeiro}\\[2pt]
Gerentes do projeto (co-gestão)

\vspace{1.5cm}

Três de Maio, 1\textsuperscript{o} de outubro de 2026
\end{titlepage}

\tableofcontents

%======================================================================
% Visão geral
%======================================================================
\capitulosemnumero{Visão Geral do Projeto}

\noindent Este documento reúne, em um só lugar, tudo o que a dupla produziu para o TrackFleet ao longo das aulas de Gestão de Projetos: um capítulo para cada um dos sete domínios de desempenho do Guia PMBOK\textregistered\ 8\textsuperscript{a} Edição, na ordem em que foram estudados. Cada capítulo aplica os conceitos da aula a decisões concretas do projeto; o último capítulo mostra como os domínios se amarram.

\secaosemnumero{O Projeto em Tópicos}
\begin{itemize}[itemsep=3pt]
  \item \textbf{O que é:} uma plataforma web de rastreamento e gestão de frotas, integrada por API a rastreadores GPS já instalados nos veículos.
  \item \textbf{Para quem:} pequenas transportadoras; o usuário é o gestor de frota.
  \item \textbf{Problema:} falta de visibilidade em tempo real da frota e de dados objetivos para decisões de manutenção e uso dos veículos, agravada pelo custo elevado das soluções corporativas.
  \item \textbf{Objetivo:} lançar a plataforma até o fim do semestre letivo.
  \item \textbf{O que entrega:} mapa quase em tempo real, quatro motores de alerta (manutenção, documentos, cerca virtual e perda de sinal), três relatórios exportáveis e um painel de indicadores.
  \item \textbf{Valor esperado:} redução de custos operacionais e de manutenção corretiva, mais segurança para veículos e motoristas e decisões baseadas em dados reais.
  \item \textbf{Patrocinador:} Diretoria da transportadora parceira (papel fictício definido pela dupla).
  \item \textbf{Gerentes do projeto:} Enzo Allebrand e Kerlon Ribeiro, em co-gestão.
  \item \textbf{Restrições:} prazo do semestre letivo, orçamento zero para hardware próprio e LGPD.
\end{itemize}

\secaosemnumero{Painel do Projeto}
\begin{center}
\small
\begin{tabularx}{\textwidth}{>{\raggedright\arraybackslash}p{4.2cm} L}
\toprule
\textbf{Indicador} & \textbf{Valor e onde está} \\
\midrule
Prazo & 31/08 a 12/11/2026, 51 dias úteis, com buffer de 24 dias úteis até o fim do semestre (17/12) --- Capítulo~\ref{cap:cronograma} \\
Caminho crítico & Requisitos, arquitetura, banco, API de cadastros, motor de alertas, relatórios, infraestrutura, integração, testes com dados reais, piloto e documentação --- Capítulo~\ref{cap:cronograma} \\
Escopo & 4 épicos, 21 atividades, 14 pacotes de trabalho na EAP --- Capítulo~\ref{cap:escopo} \\
Linha de base de custos (BAC) & R\$ 32.770, com reserva de contingência de R\$ 3.970 --- Capítulo~\ref{cap:financas} \\
Orçamento total & R\$ 34.409, com reserva de gerenciamento de R\$ 1.639 --- Capítulo~\ref{cap:financas} \\
Retorno & VPL de 3 anos de cerca de R\$ 14.588, TIR de cerca de 35\% ao ano e payback de cerca de 1,7 ano --- Capítulo~\ref{cap:financas} \\
Partes interessadas & 8 partes no registro; maior lacuna de engajamento no gestor de frota --- Capítulo~\ref{cap:partes} \\
Equipe & 2 GPs, 30 horas semanais cada, sem sobrecarga após o nivelamento --- Capítulo~\ref{cap:recursos} \\
Riscos & 8 ameaças e 3 oportunidades; maior risco R01 (API do provedor, P $\times$ I = 16) --- Capítulo~\ref{cap:riscos} \\
\bottomrule
\end{tabularx}
\end{center}

\secaosemnumero{Como Este Documento Está Organizado}
\begin{center}
\small
\begin{tabularx}{\textwidth}{l l L L}
\toprule
\textbf{Capítulo} & \textbf{Aula} & \textbf{Documentos} & \textbf{Pergunta que responde} \\
\midrule
1 Governança & 13/08 & Termo de abertura; papéis e RACI; modelo e escalonamento; métricas e sinal; regras e registro de decisões & Quem decide o quê, como e com quais critérios? \\
2 Escopo & 20/08 & Escopo do produto, do projeto e qualidade & O que entra, o que fica de fora e quando está pronto? \\
3 Cronograma & 27/08 & Atividades, dependências, estimativas, caminho crítico, linha de base & Quando cada coisa acontece, e quem faz? \\
4 Finanças & 03/09 & Plano financeiro, estimativas, orçamento, valor agregado e business case & Quanto custa e por que vale a pena? \\
5 Partes Interessadas & 10/09 & Registro; plano de engajamento; plano de comunicações & Quem influencia o projeto e como manter todos alinhados? \\
6 Recursos & 17/09 & Termo da equipe; RACI; EAR; estimativa e histograma & Quem faz o trabalho e com quais recursos? \\
7 Riscos & 24/09 & Plano de gerenciamento; registro; análise e respostas & O que pode dar errado ou certo, e o que fazer? \\
\bottomrule
\end{tabularx}
\end{center}

\noindent Dentro dos capítulos, quando um texto cita o ``Documento de Escopo'', o ``Documento de Cronograma'' e assim por diante, a referência é ao capítulo do domínio correspondente: cada capítulo é o documento que a dupla entregou na aula daquele domínio.

\secaosemnumero{Como Ler as Datas}
Todos os capítulos descrevem o plano do projeto, cuja execução começou em 31/08, na semana seguinte à aula de Cronograma. Como cada domínio foi estudado em uma aula diferente, alguns planos foram formalizados com a execução já em andamento: o Plano de Engajamento (10/09) e o registro de riscos (24/09), por exemplo, trazem ações e gatilhos das primeiras semanas, como o mapeamento da API do provedor até 04/09. Esses itens foram mantidos para que o plano fique completo e coerente do início ao fim. O acompanhamento do que de fato aconteceu é feito no check-in semanal e não faz parte deste documento; o exemplo de valor agregado do Capítulo~\ref{cap:financas} é uma simulação, não uma medição.

\secaosemnumero{Notas desta Versão}
Para o Trabalho Prático 01 todos os documentos foram revisados juntos, e as correções feitas estão registradas no log de decisões (Capítulo~\ref{cap:governanca}):
\begin{itemize}[itemsep=2pt]
  \item \textbf{Cronograma nivelado:} a primeira linha de base punha atividades do mesmo GP em paralelo e ignorava feriados; com dono único por atividade e nivelamento de recursos, o prazo passou de 49 para 51 dias úteis, absorvidos pelo buffer.
  \item \textbf{Finanças recalculadas:} os custos agora saem das atividades do cronograma nivelado, e a reserva de contingência passou a ser o VME do registro de riscos.
  \item \textbf{Escopo completado:} EAP com arquitetura e infraestrutura, dicionário com todos os pacotes, rastreabilidade dos requisitos e requisitos não funcionais com valores medíveis.
  \item \textbf{Coerência entre domínios:} RACI de Governança e de Recursos, donos do cronograma e tecnologias citadas em Partes Interessadas e Recursos passaram a usar os mesmos dados.
  \item \textbf{Domínio de Riscos:} produzido a partir da aula de 24/09.
\end{itemize}

%======================================================================
% Domínios
%======================================================================
\chapter{Governança}\label{cap:governanca}
\noindent Aula de 13/08. O primeiro domínio define quem decide o quê, com quais critérios e como o valor será medido --- a estrutura que sustenta todos os outros. O capítulo reúne as cinco partes do Documento de Governança.
\subdocumento{../governanca/termo_de_abertura}
\subdocumento{../governanca/papeis_e_raci}
\subdocumento{../governanca/modelo_e_escalonamento}
\subdocumento{../governanca/metricas_e_sinal}
\subdocumento{../governanca/regras_e_registro}

\chapter{Escopo}\label{cap:escopo}
\noindent Aula de 20/08. As três dimensões do escopo --- produto, projeto e qualidade --- tratadas separadamente, do catálogo de alertas à Definição de Pronto.

\documentclass[11pt,a4paper]{article}
\usepackage[utf8]{inputenc}
\usepackage[T1]{fontenc}
\usepackage[brazil]{babel}
\usepackage[margin=2.5cm]{geometry}
\usepackage{mathptmx}
\usepackage{amssymb}
\usepackage{enumitem}
\usepackage{booktabs}
\usepackage{tabularx}
\usepackage{array}
\usepackage{forest}
\usepackage{hyperref}
\hypersetup{colorlinks=false, pdfborder={0 0 0}}
\usepackage{fancyhdr}

\pagestyle{fancy}
\fancyhf{}
\fancyhead[L]{\small TrackFleet --- Escopo}
\fancyhead[R]{\small\thepage}
\renewcommand{\headrulewidth}{0.4pt}

\newcolumntype{L}{>{\raggedright\arraybackslash}X}
\newcommand{\dodbox}{$\square$}

\setlist{topsep=3pt}

\begin{document}

\begin{center}
{\Large\bfseries DOCUMENTO DE ESCOPO DO PROJETO}\\[2pt]
{\large\bfseries TrackFleet --- Gestão de Rastreador de Automóveis}\\[4pt]
Disciplina de Gerenciamento de Projetos $\cdot$ Framework PMBOK\textregistered\ 8\textsuperscript{a} Edição $\cdot$ Domínio: Escopo\\[2pt]
Integrantes: Enzo Allebrand e Kerlon Ribeiro (dois GPs) $\cdot$ 4\textsuperscript{o} semestre $\cdot$ Data: 20/08
\end{center}

\vspace{2mm}
\hrule
\vspace{4mm}

\noindent No PMBOK\textregistered\ 8, escopo não é uma coisa só: são três dimensões que acontecem ao mesmo tempo. Este documento reúne as três em um só lugar, cada uma em sua parte, para evitar a armadilha de comprimi-las em um único bloco:
\begin{itemize}[itemsep=2pt]
  \item \textbf{Escopo do Produto (o quê):} as características, funções e atributos do que o TrackFleet deve entregar.
  \item \textbf{Escopo do Projeto (o como):} todo o trabalho necessário, e apenas o necessário, para entregar o produto.
  \item \textbf{Qualidade:} a terceira dimensão do escopo --- não um domínio à parte, mas um atributo integrante dele.
\end{itemize}

\vspace{2mm}
\hrule
\vspace{4mm}

%======================================================================
\section{Escopo do Produto (o quê)}
%======================================================================

\noindent\textbf{\large Escopo do Produto}

\vspace{1mm}
\noindent A descrição das características, funções e atributos do resultado que o TrackFleet deve entregar. Responde à pergunta ``o que a entrega deve ser e fazer?'', distinta do Escopo do Projeto (o ``como'') e da Qualidade, tratados nas partes seguintes.

\subsection{Descrição do Produto}
\begin{itemize}[itemsep=3pt]
  \item \textbf{O que é:} uma plataforma web de rastreamento e gestão de frotas, acessada por navegador.
  \item \textbf{Para quem:} pequenas transportadoras, operadas por um gestor de frota e utilizadas também pelos motoristas.
  \item \textbf{Como funciona:} integra-se, via API, a rastreadores GPS já existentes nos veículos da transportadora parceira; não há fabricação de hardware próprio. Os dados de posição chegam do provedor, são processados pela plataforma e transformados em mapa, alertas e relatórios.
  \item \textbf{O que entrega:} visibilidade quase em tempo real da frota, alertas automáticos (manutenção, vencimento de documentos, cercas virtuais e perda de sinal) e indicadores de uso para apoiar decisões de gestão baseadas em dados reais.
\end{itemize}

\subsection{Arquitetura e Tecnologias}
Como o produto é uma plataforma de software, ``o que ele é'' inclui as decisões de tecnologia e infraestrutura que o sustentam. As escolhas abaixo são a arquitetura de referência do projeto; mudanças estruturais nela seguem o controle de mudanças definido na Parte II.

\subsubsection*{Pilha de tecnologia}
\begin{center}
\begin{tabularx}{\textwidth}{l L L}
\toprule
\textbf{Camada} & \textbf{Tecnologia} & \textbf{Por que} \\
\midrule
Backend / API & Python com Django e Django REST Framework & Framework maduro, produtivo e adequado a regras de negócio e relatórios; expõe uma API REST consumida pelo frontend. \\
Tarefas agendadas & Job agendado (comando de gerenciamento do Django via cron) & Busca periódica das posições do rastreador e verificação das regras de alerta, sem a complexidade de uma fila de mensagens dedicada. \\
Banco de dados & PostgreSQL & Armazena veículos, motoristas, trajetos e alertas; os cálculos de distância e de cerca virtual (ponto e raio) são resolvidos na própria aplicação. \\
Frontend & React (aplicação de página única) com Leaflet e Recharts & Leaflet renderiza o mapa sobre tiles do OpenStreetMap; Recharts monta os gráficos do painel; a interface é responsiva. \\
Autenticação & Tokens JWT com perfil único (gestor) & Sessões seguras; a primeira versão atende apenas o perfil de gestor. \\
Integração externa & Cliente HTTP para a API do provedor de rastreamento & Consome as posições por chamadas agendadas (polling) ou por webhook, conforme o provedor piloto suportar. \\
Geração de relatórios & Exportação em PDF e XLS & Relatórios entregues em formato imprimível (PDF) e em planilha (XLS) para análise. \\
\bottomrule
\end{tabularx}
\end{center}

\subsubsection*{Infraestrutura}
\begin{itemize}[itemsep=2pt]
  \item \textbf{Contêineres:} a aplicação roda em contêineres Docker orquestrados por docker-compose (serviços: API, banco PostgreSQL e frontend).
  \item \textbf{Borda e segurança de transporte:} proxy reverso Nginx com TLS, servindo tudo sob HTTPS.
  \item \textbf{Hospedagem:} nuvem de baixo custo ou VPS, coerente com a restrição de orçamento zero para hardware próprio.
  \item \textbf{Operação:} segredos em variáveis de ambiente, backup periódico do banco e registro de logs da aplicação.
\end{itemize}

\subsubsection*{Fluxo de dados, ponta a ponta}
\begin{enumerate}[itemsep=2pt]
  \item O rastreador GPS do veículo reporta a posição ao provedor de rastreamento de terceiros.
  \item Um job agendado consome essas posições pela API do provedor e as grava no PostgreSQL.
  \item As regras de alerta são avaliadas sobre os dados recebidos, na mesma execução periódica.
  \item O gestor visualiza mapa, alertas e relatórios no frontend React.
\end{enumerate}

\subsection{Requisitos Funcionais}
Organizados pelos módulos do produto.

\begin{itemize}[itemsep=3pt]
  \item \textbf{Cadastros:} cadastro de veículos (placa, modelo, hodômetro, documentos), de motoristas (dados, CNH e validade) e do vínculo motorista--veículo. O motorista é um registro da frota, não um usuário do sistema nesta fase.
  \item \textbf{Integração e ingestão:} integração com a API do provedor de rastreamento piloto e ingestão periódica das posições recebidas.
  \item \textbf{Monitoramento:} mapa com a localização quase em tempo real dos veículos ativos e histórico de trajetos por veículo e por período.
  \item \textbf{Cercas virtuais:} definição de áreas circulares (ponto central e raio) permitidas ou proibidas por veículo.
  \item \textbf{Motor de alertas:} geração automática dos alertas do catálogo (Seção 1.5), com registro de status (aberto, resolvido) e notificação no aplicativo e por e-mail.
  \item \textbf{Relatórios:} emissão dos relatórios do catálogo (Seção 1.6), com filtros e exportação em PDF e XLS.
  \item \textbf{Painel de indicadores:} veículos ativos e alertas pendentes por categoria.
  \item \textbf{Acesso:} autenticação com perfil de gestor.
\end{itemize}

\subsection{Catálogo de Alertas}
Detalhar ``emitir alertas'' significa dizer exatamente quais alertas, e como cada um é disparado. Em vez de uma lista longa de alertas parecidos, o TrackFleet organiza a emissão em quatro motores de alerta configuráveis, cada um cobrindo vários itens da frota. Cada alerta gerado é registrado na plataforma e notificado ao gestor no aplicativo e por e-mail.

\begin{center}
\begin{tabularx}{\textwidth}{L L L}
\toprule
\textbf{Motor de alerta} & \textbf{Itens que cobre} & \textbf{Gatilho (condição de disparo)} \\
\midrule
Manutenção preventiva (por km/tempo) & Troca de óleo, revisão periódica, pneus e freios & Quilometragem rodada ou tempo decorrido desde a última execução do item atinge o limite configurado \\
Vencimento de documento (por data) & Licenciamento/IPVA, seguro do veículo, CNH do motorista & Data de vencimento do documento se aproxima, com a antecedência configurada \\
Cerca virtual (geofence) & Áreas permitidas ou proibidas por veículo & O veículo cruza o raio de uma área definida (ponto central e raio) \\
Perda de sinal do rastreador & Todos os veículos ativos & Nenhuma posição recebida do veículo por tempo acima do limite \\
\bottomrule
\end{tabularx}
\end{center}

\noindent Os limites (quilometragem, antecedência de vencimento, ponto e raio da cerca, tempo sem sinal) e os itens de cada motor são parâmetros configuráveis pelo gestor, não valores fixos no código. Alertas que dependeriam de telemetria não garantida pela API do provedor (por exemplo, excesso de velocidade e ociosidade de motor) ficam fora desta fase, conforme a Seção 1.7.

\subsection{Catálogo de Relatórios}
Detalhar ``emitir relatórios'' significa dizer quais relatórios e o que cada um responde. Todos são filtráveis por período, veículo e motorista, e exportáveis em PDF e XLS.

\begin{center}
\begin{tabularx}{\textwidth}{L L}
\toprule
\textbf{Relatório} & \textbf{O que responde} \\
\midrule
Relatório de alertas & Quais alertas foram emitidos no período, por categoria, veículo e motorista, com o status de cada um (aberto ou resolvido); inclui os documentos próximos do vencimento \\
Relatório de uso por veículo & Km rodados, horas em operação e número de viagens de cada veículo \\
Relatório de manutenção & Manutenções previstas frente às realizadas por veículo \\
\bottomrule
\end{tabularx}
\end{center}

\noindent O painel gerencial (dashboard), com veículos ativos e alertas pendentes, é uma visão de tela do produto (Seção 1.3), não um relatório exportável.

\subsection{Requisitos Não Funcionais}
\begin{itemize}[itemsep=2pt]
  \item \textbf{Conformidade com a LGPD:} consentimento informado e minimização no tratamento de dados de localização e de motoristas; retenção do histórico por prazo limitado; acesso restrito por autenticação.
  \item \textbf{Atualidade da posição:} a localização exibida no mapa reflete a última leitura do rastreador com atraso perceptível mínimo, respeitado o intervalo de envio do provedor.
  \item \textbf{Segurança:} tráfego sob HTTPS, senhas armazenadas com hash e controle de acesso autenticado.
  \item \textbf{Responsividade:} interface utilizável em desktop e em navegador mobile.
  \item \textbf{Disponibilidade:} plataforma no ar durante toda a janela de testes com a transportadora piloto.
  \item \textbf{Desempenho:} telas principais (mapa e painel) respondem em tempo adequado ao uso do gestor.
\end{itemize}

\subsection{Fora do Escopo do Produto (Exclusões)}
Tão importante quanto listar o que entra é deixar explícito o que não entra, para evitar expectativas equivocadas:
\begin{itemize}[itemsep=2pt]
  \item Fabricação, aquisição ou instalação física de rastreadores GPS próprios.
  \item Aplicativo mobile nativo (iOS/Android); a primeira versão é web responsiva.
  \item Módulo de pagamento, faturamento ou cobrança.
  \item Roteirização ou otimização automática de rotas.
  \item Integração simultânea com múltiplos provedores de rastreamento (apenas um provedor piloto nesta fase).
  \item Bloqueio remoto do veículo ou qualquer atuação sobre o hardware embarcado.
  \item Login e aplicativo próprios para o motorista; nesta fase o motorista é apenas um registro da frota, e o único usuário do sistema é o gestor.
  \item Alertas que dependem de telemetria não garantida pela API do provedor: excesso de velocidade, ociosidade de motor (motor ligado parado) e uso fora do horário de expediente.
  \item Cercas virtuais desenhadas como polígonos livres; nesta fase a cerca é circular (ponto central e raio).
  \item Indicador de economia estimada em reais, por depender de premissas de custo frágeis nesta fase.
\end{itemize}

\subsection{Estrutura do Escopo: Backlog do Produto}
Por se tratar de um projeto conduzido com governança híbrida e leve (conforme o Documento de Governança), o escopo do produto é estruturado como um backlog, decomposto em épicos, features e histórias de usuário, em vez de uma EAP rígida.

\begin{itemize}[itemsep=4pt]
  \item[\textbf{E1}] Monitoramento em tempo real --- localização atual e histórico de trajetos.
  \begin{itemize}
    \item Como gestor, quero ver todos os veículos ativos em um mapa, para saber onde a frota está agora.
    \item Como gestor, quero consultar o trajeto de um veículo em um período, para auditar o uso.
  \end{itemize}
  \item[\textbf{E2}] Alertas automáticos --- manutenção, vencimento de documentos, cercas virtuais e perda de sinal.
  \begin{itemize}
    \item Como gestor, quero receber um alerta antes que a manutenção vire corretiva, para reduzir custo.
    \item Como gestor, quero ser avisado quando um documento estiver perto de vencer, para não rodar irregular.
    \item Como gestor, quero um alerta quando um veículo sair de uma área permitida, para agir na hora.
  \end{itemize}
  \item[\textbf{E3}] Gestão de frota e usuários --- cadastros, vínculos e perfis.
  \begin{itemize}
    \item Como gestor, quero cadastrar veículos e motoristas e vinculá-los, para organizar a frota.
  \end{itemize}
  \item[\textbf{E4}] Relatórios e indicadores --- painel e exportação.
  \begin{itemize}
    \item Como gestor, quero um painel com veículos ativos e alertas pendentes, para decidir rapidamente.
    \item Como gestor, quero exportar o relatório de alertas do período, para prestar contas à diretoria.
  \end{itemize}
\end{itemize}

\subsection{Value Breakdown Structure (VBS)}
A VBS conecta o escopo do produto ao valor pretendido, tornando visível qual entregável move mais o ponteiro do projeto.

\begin{center}
\begin{tabularx}{\textwidth}{L r}
\toprule
\textbf{Entregável (épico)} & \textbf{Valor} \\
\midrule
E1 --- Monitoramento em tempo real & 35\% \\
E2 --- Alertas automáticos & 35\% \\
E4 --- Relatórios e indicadores & 20\% \\
E3 --- Gestão de frota e usuários (infraestrutura base) & 10\% \\
Conformidade com a LGPD & 100\% (obrigatório) \\
\bottomrule
\end{tabularx}
\end{center}

\noindent Leitura da VBS: o monitoramento em tempo real e os alertas automáticos são os entregáveis mais valiosos e devem ser priorizados sobre relatórios avançados caso o prazo do semestre aperte. A conformidade com a LGPD é obrigatória e não é negociável, independentemente de prioridade de valor.

\vspace{2mm}
\hrule
\vspace{4mm}

%======================================================================
\section{Escopo do Projeto (o como)}
%======================================================================

\noindent\textbf{\large Escopo do Projeto}

\vspace{1mm}
\noindent Todo o trabalho necessário, e apenas o necessário, para entregar o Escopo do Produto descrito na Parte I. É o componente mais importante da linha de base, pois nele reside o valor esperado do projeto.

\subsection{Descrição do Trabalho do Projeto}
\begin{itemize}[itemsep=2pt]
  \item Levantar requisitos com a transportadora parceira e mapear a integração com a API do provedor de rastreamento piloto.
  \item Projetar a arquitetura da solução (backend Django, frontend React, banco PostgreSQL, integração externa e infraestrutura em contêineres).
  \item Desenvolver o backend: API própria, ingestão das posições e os quatro motores de alerta do catálogo.
  \item Desenvolver o frontend: mapa, painel de indicadores, cadastros e relatórios exportáveis.
  \item Implementar a autenticação do gestor e as medidas de conformidade com a LGPD.
  \item Testar a solução (testes automatizados e testes com dados reais/simulados do rastreador).
  \item Homologar a plataforma com a transportadora piloto.
  \item Treinar o gestor piloto e coletar feedback.
  \item Documentar o projeto e encerrar a primeira fase, com registro de lições aprendidas.
\end{itemize}

\subsection{Estrutura Analítica do Projeto (EAP)}
A EAP decompõe o trabalho total em pacotes menores, atribuíveis e mensuráveis, tornando o acordo sobre o escopo visível a todos os envolvidos. O projeto está no topo, os grandes entregáveis no segundo nível e os pacotes de trabalho no terceiro.

\begin{center}
\resizebox{\textwidth}{!}{%
\begin{forest}
  for tree={
    draw,
    rounded corners,
    align=center,
    inner sep=5pt,
    l sep=16pt,
    s sep=8pt,
    edge={-},
    font=\small,
  }
  [{1\\TrackFleet}
    [{1.1\\Gestão do\\Projeto}
      [{1.1.1\\Termo de abertura\\e governança}]
      [{1.1.2\\Planejamento e\\acompanhamento}]
    ]
    [{1.2\\Requisitos}
      [{1.2.1\\Entrevistas com a\\transportadora}]
      [{1.2.2\\Especificação de\\requisitos}]
    ]
    [{1.3\\Desenvolvimento}
      [{1.3.1\\Backend, API e\\ingestão}]
      [{1.3.2\\Motor de\\alertas}]
      [{1.3.3\\Frontend (mapa,\\painel, relatórios)}]
      [{1.3.4\\Autenticação (gestor)\\e LGPD}]
    ]
    [{1.4\\Testes e\\Homologação}
      [{1.4.1\\Testes\\automatizados}]
      [{1.4.2\\Piloto com a\\transportadora}]
    ]
    [{1.5\\Treinamento e\\Encerramento}
      [{1.5.1\\Treinamento do\\gestor piloto}]
      [{1.5.2\\Documentação e\\lições aprendidas}]
    ]
  ]
\end{forest}
}
\end{center}

\subsection{Dicionário da EAP}
\begin{center}
\begin{tabularx}{\textwidth}{l L L}
\toprule
\textbf{Pacote} & \textbf{Descrição} & \textbf{Critério de aceite} \\
\midrule
1.3.1 Backend, API e ingestão & API própria em Django consumindo dados do rastreador piloto e persistindo posições no PostgreSQL & Dados do rastreador chegam à plataforma e ficam disponíveis para consulta \\
1.3.2 Motor de alertas & Os quatro motores que avaliam os dados e geram os alertas do catálogo, com status e notificação & Cada motor do catálogo dispara corretamente na condição configurada \\
1.3.3 Frontend & Mapa, painel, cadastros e relatórios exportáveis em React & Gestor localiza um veículo e visualiza alertas pendentes em poucos cliques \\
1.3.4 Autenticação e LGPD & Perfil de gestor e medidas de minimização e proteção de dados & Acesso autenticado e sem exposição indevida de dados pessoais \\
1.4.2 Piloto & Uso real da plataforma pela transportadora parceira por um período de teste & Transportadora usa a plataforma e valida os indicadores gerados \\
\bottomrule
\end{tabularx}
\end{center}

\subsection{Fora do Escopo do Projeto (Exclusões)}
\begin{itemize}[itemsep=2pt]
  \item Compra, instalação física ou manutenção de rastreadores GPS.
  \item Suporte comercial e operação em produção após o encerramento do piloto acadêmico.
  \item Integração com mais de um provedor de hardware de rastreamento.
  \item Certificações formais de segurança ou auditorias externas de compliance além da adequação básica à LGPD.
  \item Desenvolvimento de aplicativo mobile nativo.
\end{itemize}

\subsection{Linha de Base do Escopo e Controle de Mudanças}
A linha de base do escopo deste documento é aprovada pelo patrocinador (Diretoria da transportadora parceira), conforme definido no Documento de Governança. Qualquer mudança de escopo que altere objetivo, prazo ou entregáveis essenciais segue o mesmo fluxo de escalonamento já definido: sem consenso entre os dois GPs, ou se a mudança afetar objetivo/prazo/valor esperado, a decisão sobe ao patrocinador e é registrada no log de decisões do projeto. Isso evita o ``scope creep'', o crescimento não controlado do escopo por pequenos acréscimos aceitos por fora do processo.

\subsection{Indicadores de Escopo}
\begin{center}
\begin{tabularx}{\textwidth}{L l l}
\toprule
\textbf{Indicador} & \textbf{Fórmula} & \textbf{Leitura} \\
\midrule
Precisão da Definição & Itens entregues corretamente $\div$ Itens planejados $\times$ 100 & Maior é melhor \\
Scope Creep & Entregas não planejadas $\div$ Total de entregas $\times$ 100 & Menor é melhor \\
Estabilidade de Requisitos & Requisitos não alterados $\div$ Total de requisitos $\times$ 100 & Maior é melhor \\
\bottomrule
\end{tabularx}
\end{center}

\noindent Esses indicadores devem ser acompanhados a cada check-in do projeto (conforme a cadência definida na Governança), para que desvios de escopo sejam percebidos antes da entrega final, e não durante ela.

\vspace{2mm}
\hrule
\vspace{4mm}

%======================================================================
\section{Qualidade}
%======================================================================

\noindent\textbf{\large Qualidade}

\vspace{1mm}
\noindent A terceira dimensão do escopo: não um domínio à parte, mas um atributo integrante dele. A qualidade foca em atender às expectativas dos stakeholders e em cumprir os requisitos definidos tanto no Escopo do Produto quanto no Escopo do Projeto.

\subsection{Requisitos de Qualidade do Produto}
\begin{itemize}[itemsep=2pt]
  \item Precisão e atualidade da localização exibida no mapa, com latência aceitável em relação à posição real do rastreador.
  \item Confiabilidade dos alertas, com o mínimo possível de falsos positivos e falsos negativos em cada motor do catálogo.
  \item Usabilidade do painel: o gestor deve conseguir localizar um veículo e visualizar alertas pendentes em poucos cliques.
  \item Conformidade com a LGPD no tratamento de dados de localização e de motoristas (consentimento informado e minimização de dados coletados).
  \item Disponibilidade da plataforma durante toda a janela de testes com a transportadora piloto.
\end{itemize}

\subsection{Requisitos de Qualidade do Projeto (Processo)}
\begin{itemize}[itemsep=2pt]
  \item Todo código é revisado pelos dois GPs antes de ser integrado --- não há code review unilateral.
  \item Rotinas críticas (integração com a API do rastreador, ingestão e cálculo de alertas) possuem testes automatizados.
  \item A documentação do projeto é atualizada a cada entrega, não deixada para o final.
\end{itemize}

\subsection{Definição de Pronto (DoD)}
A DoD é o conjunto de critérios que uma funcionalidade precisa cumprir para ser considerada apta ao uso pela transportadora piloto. Ela é definida antes de a funcionalidade começar a ser construída, para blindar o aceite e evitar a conversa em que a equipe diz que terminou e o patrocinador diz que não.

\begin{itemize}[itemsep=3pt]
  \item[\dodbox] Funcionalidade implementada conforme o requisito descrito no Escopo do Produto.
  \item[\dodbox] Código revisado por ambos os GPs (Enzo Allebrand e Kerlon Ribeiro).
  \item[\dodbox] Testada com dados reais ou simulados do rastreador GPS piloto.
  \item[\dodbox] Sem exposição indevida de dados pessoais (checagem de conformidade com a LGPD).
  \item[\dodbox] Validada e aceita pelos dois GPs; mudanças de escopo, prazo ou objetivo passam ainda pelo patrocinador, conforme a matriz RACI da Governança.
\end{itemize}

\subsection{Critérios de Aceite por Entregável}
\begin{center}
\begin{tabularx}{\textwidth}{L L}
\toprule
\textbf{Entregável} & \textbf{Critério de aceite} \\
\midrule
Monitoramento em tempo real (E1) & O gestor visualiza, no mapa, a posição atual de todos os veículos ativos com atraso perceptível mínimo \\
Alertas automáticos (E2) & Cada motor de alerta do catálogo dispara automaticamente na condição configurada, sem falso positivo ou negativo relevante \\
Relatórios e indicadores (E4) & O painel exibe corretamente veículos ativos e alertas pendentes, e os relatórios são exportáveis em PDF e XLS \\
Piloto com a transportadora parceira & A transportadora utiliza a plataforma no período de teste e confirma, por escrito, que os indicadores refletem a realidade da frota \\
\bottomrule
\end{tabularx}
\end{center}

\subsection{Validar \texttimes\ Controlar o Escopo}
As duas atividades são complementares, mas distintas:
\begin{itemize}[itemsep=3pt]
  \item \textbf{Validar o Escopo} --- aceite formal, externo: o patrocinador (Diretoria da transportadora parceira) aprova cada marco relevante do projeto.
  \item \textbf{Controlar o Escopo} --- verificação interna, contínua: os dois GPs conferem, a cada check-in, se o que está sendo construído continua dentro da linha de base aprovada.
\end{itemize}
\noindent Um valida o resultado; o outro cuida do processo. Confundir os dois é um dos erros mais comuns de gestão de escopo.

\subsection{Prevenção de Scope Creep}
As regras de proteção (guardrails) já definidas no Documento de Governança sustentam também a qualidade do escopo:
\begin{itemize}[itemsep=2pt]
  \item Toda mudança relevante de escopo, prazo ou objetivo passa pelo patrocinador antes de ser aceita.
  \item Toda decisão importante é registrada no log de decisões do projeto, no momento em que é tomada.
  \item Nenhuma funcionalidade é considerada ``pronta'' sem passar pela DoD definida na Seção 3.3.
  \item Sem consenso entre os dois GPs, a decisão sobe direto ao patrocinador, que desempata.
\end{itemize}

\subsection{Métricas de Qualidade}
\begin{center}
\begin{tabularx}{\textwidth}{L l}
\toprule
\textbf{Métrica} & \textbf{Leitura} \\
\midrule
Taxa de defeitos encontrados no piloto & Menor é melhor \\
\% de funcionalidades aceitas sem retrabalho na primeira validação & Maior é melhor \\
\% de alertas corretos (sem falso positivo/negativo) & Maior é melhor \\
\bottomrule
\end{tabularx}
\end{center}

\end{document}


\chapter{Cronograma}\label{cap:cronograma}
\noindent Aula de 27/08, revisado para o Trabalho Prático 01. Da EAP às atividades, dependências e estimativas, até o caminho crítico com recursos, a linha de base e o buffer.

\documentclass[11pt,a4paper]{article}
\usepackage[utf8]{inputenc}
\usepackage[T1]{fontenc}
\usepackage[brazil]{babel}
\usepackage[margin=2cm]{geometry}
\usepackage{mathptmx}
\usepackage{amssymb}
\usepackage{amsmath}
\usepackage{enumitem}
\usepackage{booktabs}
\usepackage{tabularx}
\usepackage{array}
\usepackage{tikz}
\usetikzlibrary{arrows.meta,positioning}
\usepackage{hyperref}
\hypersetup{colorlinks=false, pdfborder={0 0 0}}
\usepackage{fancyhdr}

\pagestyle{fancy}
\fancyhf{}
\fancyhead[L]{\small TrackFleet --- Cronograma}
\fancyhead[R]{\small\thepage}
\renewcommand{\headrulewidth}{0.4pt}

\newcolumntype{L}{>{\raggedright\arraybackslash}X}
\newcolumntype{C}{>{\centering\arraybackslash}X}
\setlist{topsep=3pt}

\begin{document}

\begin{center}
{\Large\bfseries DOCUMENTO DE CRONOGRAMA DO PROJETO}\\[2pt]
{\large\bfseries TrackFleet --- Gestão de Rastreador de Automóveis}\\[4pt]
Disciplina de Gerenciamento de Projetos $\cdot$ Framework PMBOK\textregistered\ 8\textsuperscript{a} Edição $\cdot$ Domínio: Cronograma\\[2pt]
Integrantes: Enzo Allebrand e Kerlon Ribeiro (dois GPs) $\cdot$ 4\textsuperscript{o} semestre $\cdot$ Data: 27/08
\end{center}

\vspace{2mm}
\hrule
\vspace{4mm}

\noindent No PMBOK\textregistered\ 8, o domínio de desempenho Cronograma cobre três processos: \textbf{Definir Atividades} (quebrar a EAP em trabalho executável), \textbf{Sequenciar Atividades} (ligar esse trabalho por dependências lógicas) e \textbf{Desenvolver e Controlar o Cronograma} (transformar a rede de atividades em datas, identificar o caminho crítico e manter isso vivo até a entrega). Como a aula ensina, o cronograma do TrackFleet não é a barra de Gantt em si: é, ao mesmo tempo, um modelo de como o projeto se comporta se algo mudar, uma comunicação entre os dois GPs e o patrocinador sobre quem faz o quê e quando, e um compromisso --- a linha de base contra a qual o realizado é medido.

\vspace{2mm}
\hrule
\vspace{4mm}

%======================================================================
\section{Definir Atividades: da EAP ao trabalho executável}
%======================================================================

\noindent A Estrutura Analítica do Projeto (EAP) do TrackFleet já foi decomposta em pacotes de trabalho no Documento de Escopo (Seção 2.2: 1.1 Gestão do Projeto, 1.2 Requisitos, 1.3 Desenvolvimento, 1.4 Testes e Homologação, 1.5 Treinamento e Encerramento). Aqui, cada pacote relevante à execução é quebrado em atividades no padrão ensinado na aula: \textbf{regra 8/80} (entre 8 e 80 horas de esforço, ou seja, de 1 a 10 dias úteis de trabalho de uma pessoa), \textbf{verbo no início} (ação clara, não um tema genérico) e \textbf{um dono só} por atividade.

\begin{center}
\footnotesize
\begin{tabularx}{\textwidth}{l L l l c}
\toprule
\textbf{ID} & \textbf{Atividade} & \textbf{Pacote EAP} & \textbf{Responsável} & \textbf{Duração} \\
\midrule
A1  & Levantar requisitos com a transportadora parceira & 1.2.1 & Ambos            & 3 dias \\
A2  & Mapear a API do provedor de rastreamento piloto & 1.2.1 & Kerlon           & 2 dias \\
A3  & Especificar requisitos funcionais e não funcionais & 1.2.2 & Enzo             & 3 dias \\
A4  & Projetar a arquitetura da solução & 1.1.2 & Ambos            & 2 dias \\
A5  & Modelar o banco de dados (veículos, motoristas, posições, alertas) & 1.3.1 & Enzo             & 3 dias \\
A6  & Implementar a API de cadastros e vínculos & 1.3.1 & Enzo             & 6 dias \\
A7  & Implementar a ingestão periódica das posições do provedor & 1.3.1 & Kerlon           & 6 dias \\
A8  & Implementar os quatro motores do catálogo de alertas & 1.3.2 & Enzo             & 9 dias \\
A9  & Implementar a autenticação JWT do gestor & 1.3.4 & Kerlon           & 2 dias \\
A10 & Aplicar as medidas de conformidade com a LGPD & 1.3.4 & Kerlon           & 3 dias \\
A11 & Construir o mapa (Leaflet) com localização quase em tempo real & 1.3.3 & Kerlon           & 6 dias \\
A12 & Construir o painel de indicadores (Recharts) & 1.3.3 & Kerlon           & 4 dias \\
A13 & Construir as telas de cadastro de veículos e motoristas & 1.3.3 & Kerlon           & 3 dias \\
A14 & Construir os relatórios exportáveis (PDF/XLS) & 1.3.3 & Kerlon           & 4 dias \\
A15 & Escrever testes automatizados das rotinas críticas & 1.4.1 & Ambos            & 4 dias \\
A16 & Integrar backend e frontend em ambiente de homologação & 1.4.1 & Ambos            & 3 dias \\
A17 & Testar com dados reais/simulados do rastreador piloto & 1.4.1 & Ambos            & 3 dias \\
A18 & Homologar a plataforma com a transportadora piloto & 1.4.2 & Ambos            & 11 dias \\
A19 & Treinar o gestor piloto & 1.5.1 & Ambos            & 1 dia \\
A20 & Documentar o projeto e registrar lições aprendidas & 1.5.2 & Ambos            & 3 dias \\
\bottomrule
\end{tabularx}
\end{center}

\noindent A atividade A18 (piloto) foge da faixa 8/80 de propósito: não é um pacote de esforço contínuo da equipe, e sim uma janela de operação assistida, cuja duração é ditada pelo tempo mínimo de uso real necessário para a transportadora validar os indicadores gerados --- não pelo número de horas que os GPs trabalham nela.

\vspace{2mm}
\hrule
\vspace{4mm}

%======================================================================
\section{Sequenciar Atividades: dependências}
%======================================================================

\noindent O Método de Diagrama de Precedência (PDM) liga as atividades por quatro tipos de relação. O TrackFleet usa majoritariamente Término$\to$Início (FS), com alguns pares Início$\to$Início (SS) deliberados para ganhar tempo; não há nenhuma dependência Início$\to$Término (SF) no cronograma, coerente com a raridade desse tipo apontada na aula.

\begin{center}
\footnotesize
\begin{tabularx}{\textwidth}{l L l L}
\toprule
\textbf{ID} & \textbf{Predecessora(s)} & \textbf{Tipo} & \textbf{Natureza e observação} \\
\midrule
A3  & A1                  & FS & Obrigatória, interna --- não dá para especificar sem ter levantado o requisito. \\
A4  & A3                  & FS & Obrigatória, interna. \\
A5  & A4                  & FS & Obrigatória, interna. \\
A9  & A4                  & SS & Arbitrária, interna --- autenticação começa junto com a modelagem de dados, sem esperar o backend inteiro. \\
A6  & A5                  & FS & Obrigatória, interna --- precisa do modelo de dados fechado. \\
A7  & A5, A2              & FS & A5 é obrigatória interna; A2 é \textbf{obrigatória externa} --- a ingestão depende do acesso e da documentação da API do provedor de rastreamento, fora do controle da dupla. \\
A10 & A9                  & FS & Obrigatória, interna. \\
A8  & A6, A7               & FS & Obrigatória, interna --- o motor de alertas lê cadastros e posições já ingeridas. \\
A11 & A6, A7               & FS & Obrigatória, interna --- o mapa depende da API de cadastro e da ingestão de posições. \\
A13 & A6                  & FS & Obrigatória, interna. \\
A12 & A8                  & FS & Obrigatória, interna --- o painel exibe alertas pendentes. \\
A14 & A8                  & FS & Obrigatória, interna --- o relatório de alertas depende do motor pronto. \\
A15 & A8                  & SS (lead de 6 dias) & Arbitrária, interna --- os testes automatizados começam antes do motor de alertas fechar por completo, para não empurrar tudo para o fim. \\
A16 & A9, A10, A11, A12, A13, A14 & FS & Obrigatória, interna --- integração só faz sentido com os módulos prontos. \\
A17 & A16, A15             & FS & Obrigatória, interna. \\
A18 & A17                  & FS (lag de 2 dias) & Obrigatória, interna --- os 2 dias de lag são o prazo para o patrocinador formalizar o aceite antes de a transportadora começar a usar a plataforma de verdade. \\
A19 & A18                  & SS & Arbitrária, interna --- o treinamento do gestor começa junto com o piloto. \\
A20 & A18 (FF), A19 (FS)    & FF / FS & Obrigatória, interna --- a documentação e as lições aprendidas fecham exatamente junto com o encerramento do piloto. \\
\bottomrule
\end{tabularx}
\end{center}

\vspace{2mm}
\hrule
\vspace{4mm}

%======================================================================
\section{Estimar Durações: técnica por tipo de atividade}
%======================================================================

\noindent Nem toda atividade é estimada da mesma forma. O TrackFleet combina as quatro técnicas da aula conforme a natureza do trabalho:

\begin{itemize}[itemsep=3pt]
  \item \textbf{Especialista:} atividades pequenas e bem conhecidas da dupla (A9 --- autenticação JWT, A19 --- treinamento) são estimadas por julgamento direto.
  \item \textbf{Análoga:} A1 (levantamento de requisitos) e A20 (documentação) são estimadas comparando com entregas anteriores parecidas da própria disciplina (Termo de Abertura e Documento de Escopo).
  \item \textbf{Paramétrica:} A13 (telas de cadastro) e A14 (relatórios) usam uma métrica simples --- número de telas/relatórios do catálogo (Seção 1.6 do Documento de Escopo) multiplicado pelo tempo médio por unidade.
  \item \textbf{Três pontos (PERT):} reservada para as cinco atividades mais incertas ou mais críticas do projeto, detalhadas abaixo.
\end{itemize}

\subsection{Estimativa de três pontos das atividades críticas}
A fórmula da distribuição Beta pondera o cenário mais provável e usa a distância entre o otimista e o pessimista como medida de risco:
\[
t_E = \frac{O + 4M + P}{6} \qquad \text{Desvio} \approx \frac{P - O}{6}
\]

\begin{center}
\begin{tabularx}{\textwidth}{L c c c c c}
\toprule
\textbf{Atividade} & \textbf{O} & \textbf{M} & \textbf{P} & $t_E$ & \textbf{Desvio} \\
\midrule
A6 --- API de cadastros e vínculos & 3 & 5 & 9 & 5,3 dias & 1,0 \\
A7 --- Ingestão de posições do provedor & 3 & 5 & 10 & 5,5 dias & 1,2 \\
A8 --- Motor de alertas (4 motores) & 5 & 8 & 15 & 8,7 dias & 1,7 \\
A11 --- Mapa com localização quase em tempo real & 3 & 5 & 9 & 5,3 dias & 1,0 \\
A18 --- Piloto com a transportadora parceira & 7 & 10 & 15 & 10,3 dias & 1,3 \\
\bottomrule
\end{tabularx}
\end{center}

\noindent A leitura desses números: o motor de alertas (A8) tem o maior desvio (1,7 dia) porque reúne quatro lógicas de disparo diferentes (Seção 1.5 do Documento de Escopo) na mesma atividade --- é o maior risco de estouro do cronograma, e por isso concentra a atenção das solução-chave da Seção 6. Os valores de $t_E$, arredondados para cima, são os que aparecem na tabela de atividades da Seção 1 (A6=6, A7=6, A8=9, A11=6, A18=11 dias) e alimentam o caminho crítico da Seção 4.

\vspace{2mm}
\hrule
\vspace{4mm}

%======================================================================
\section{Desenvolver o Cronograma: caminho crítico}
%======================================================================

\noindent O caminho crítico é a sequência mais longa de atividades dependentes: define a menor duração possível do projeto e é onde qualquer atraso vira atraso do projeto inteiro. Somando as durações da Seção 1 segundo as dependências da Seção 2 (dias úteis, contados a partir do início da execução):

\begin{center}
\resizebox{\textwidth}{!}{%
\begin{tikzpicture}[
  node distance=4mm and 5mm,
  box/.style={draw, rounded corners, align=center, font=\scriptsize, inner sep=3pt, minimum height=9mm, text width=19mm},
  crit/.style={box, very thick},
  arr/.style={-{Latex[length=2mm]}}
]
\node[crit] (a) {A1+A3+A4+A5\\Requisitos,\\arquitetura\\e modelagem\\11 dias};
\node[crit, right=of a] (b) {A6\\Cadastros\\API\\6 dias};
\node[crit, right=of b] (d) {A8\\Motor de\\alertas\\9 dias};
\node[crit, right=of d] (e) {A12/A14\\Painel e\\relatórios\\4 dias};
\node[crit, right=of e] (f) {A16\\Integração\\3 dias};
\node[crit, right=of f] (g) {A17\\Testes com\\dados reais\\3 dias};
\node[crit, right=of g] (h) {A18\\Piloto\\(+2d lag)\\11 dias};
\node[box, below=of d] (c) {A7\\Ingestão\\6 dias};
\node[box, below right=8mm and 5mm of c] (m) {A11/A13\\Mapa e\\cadastro\\6 dias};

\draw[arr] (a) -- (b);
\draw[arr] (a.south) -- ++(0,-4mm) -| (c.west);
\draw[arr] (b) -- (d);
\draw[arr] (d) -- (e);
\draw[arr] (e) -- (f);
\draw[arr] (f) -- (g);
\draw[arr] (g) -- (h);
\draw[arr] (b.south) |- (m.west);
\draw[arr] (c) -- (m);
\draw[arr] (m.east) -| (f.south);
\end{tikzpicture}
}
\end{center}

\noindent\textbf{Caminho crítico:} A1$\to$A3$\to$A4$\to$A5$\to$A6$\to$A8$\to$(A12 ou A14)$\to$A16$\to$A17$\to$A18 = \textbf{49 dias úteis} ($\approx$ 10 semanas). A7 corre em paralelo a A6 e chega no mesmo instante (dia 17) --- as duas são igualmente críticas, sem folga.

\subsection{Folga (float)}
\begin{itemize}[itemsep=2pt]
  \item \textbf{Folga total do ramo Mapa/Cadastro (A11/A13):} termina no dia 23, mas só é consumido pela integração (A16) no dia 30 --- \textbf{7 dias de folga total}. Pode atrasar até esse limite sem atrasar o projeto.
  \item \textbf{Folga total do ramo Autenticação/LGPD (A9/A10):} termina no dia 13, também só consumido no dia 30 --- \textbf{17 dias de folga}. É a maior folga do cronograma e, por isso, a mais arriscada: como ensina a aula, ``toda folga é gasta'' --- por isso o check-in semanal (Seção 6) cobra status binário dessa frente mesmo com folga grande, para não deixá-la ser consumida pela Lei de Parkinson.
  \item \textbf{Folga zero (caminho crítico):} A1, A3, A4, A5, A6, A7, A8, A12, A14, A16, A17 e A18 não têm respiro --- atrasar qualquer uma delas em 1 dia atrasa o projeto em 1 dia.
\end{itemize}

\subsection{Nivelamento de recursos no caminho crítico}
A6 (cadastros) e A7 (ingestão) caem em paralelo, ambas no caminho crítico, terminando no mesmo dia. Se as duas fossem atribuídas à mesma pessoa, ela seria o ``único especialista em todo caminho crítico'' --- exatamente a pegadinha de recurso superlotado da aula. Por isso A6 é do Enzo e A7 é do Kerlon: o paralelismo do diagrama só é real se os dois donos forem pessoas diferentes. Já A11 e A13, ambas do Kerlon, cabem dentro da folga de 7 dias do ramo Mapa/Cadastro mesmo sendo feitas em sequência por uma única pessoa --- não há necessidade de dividir esse par entre os dois GPs.

\vspace{2mm}
\hrule
\vspace{4mm}

%======================================================================
\section{Linha de Base, Calendário e Buffer do Projeto}
%======================================================================

\subsection{Calendário por semana}
Execução iniciando na semana de 31/08 (semana seguinte a esta aula), em blocos de 5 dias úteis:

\begin{center}
\scriptsize
\begin{tabularx}{\textwidth}{p{3.5cm}*{10}{C}}
\toprule
\textbf{Pacote} & \textbf{S1} & \textbf{S2} & \textbf{S3} & \textbf{S4} & \textbf{S5} & \textbf{S6} & \textbf{S7} & \textbf{S8} & \textbf{S9} & \textbf{S10} \\
\midrule
Requisitos e arquitetura (A1--A5) & X & X & X &  &  &  &  &  &  &  \\
Autenticação e LGPD (A9, A10) & & X & X &  &  &  &  &  &  &  \\
Backend: cadastros $\parallel$ ingestão (A6, A7) & & & X & X &  &  &  &  &  &  \\
Frontend: mapa e cadastro (A11, A13) & & & & X & X &  &  &  &  &  \\
Motor de alertas (A8) & & & & X & X & X &  &  &  &  \\
Testes automatizados (A15) & & & & & X & X &  &  &  &  \\
Frontend: painel e relatórios (A12, A14) & & & & & & X &  &  &  &  \\
Integração backend/frontend (A16) & & & & & & & X &  &  &  \\
Testes com dados reais (A17) & & & & & & & X & X &  &  \\
Piloto com a transportadora (A18) & & & & & & & & X & X & X \\
Treinamento do gestor piloto (A19) & & & & & & & & X &  &  \\
Documentação e encerramento (A20) & & & & & & & &  &  & X \\
\bottomrule
\end{tabularx}
\end{center}

\noindent A semana 10 encerra por volta de 02/11. O objetivo do termo de abertura é lançar a plataforma ``até o fim do semestre letivo'' (dezembro); da semana 10 até o encerramento do semestre sobram cerca de 5 a 6 semanas.

\subsection{Buffer do projeto (Corrente Crítica)}
Seguindo a solução-chave nº1 da aula, o TrackFleet não infla cada atividade individualmente com uma gordura ``por segurança''. As estimativas da Seção 3 já são dadas em faixa (três pontos), sem colchão escondido, e a margem de 5 a 6 semanas entre o fim da Semana 10 e o fim do semestre é o \textbf{buffer único e visível do projeto} --- agregado no final, não escondido em cada tarefa. Consumir esse buffer não é ``atraso'': é usar exatamente a reserva planejada. Só é motivo de escalonamento ao patrocinador (conforme o Documento de Governança) se o próprio buffer estiver perto de se esgotar.

\subsection{Linha de base e re-baseline}
A linha de base do cronograma --- as datas deste documento --- é aprovada pelo patrocinador (Diretoria da transportadora parceira), na mesma lógica de aprovação da linha de base de escopo (Documento de Escopo, Seção 2.6). Um novo baseline só é feito após uma mudança formal de escopo ou prazo escalada ao patrocinador; nunca para ``esconder'' um atraso que ainda cabe dentro do buffer.

\vspace{2mm}
\hrule
\vspace{4mm}

%======================================================================
\section{Controlar o Cronograma}
%======================================================================

\subsection{Validar $\times$ Controlar}
\begin{itemize}[itemsep=3pt]
  \item \textbf{Validar o cronograma} --- aceite formal, externo: o patrocinador aprova a linha de base e qualquer re-baseline.
  \item \textbf{Controlar o cronograma} --- verificação interna, contínua: os dois GPs comparam o realizado contra a linha de base a cada check-in (10 minutos, cadência definida no Documento de Governança) e registram o status binário (pronto/não pronto) de cada atividade em andamento.
\end{itemize}

\subsection{Indicadores de Cronograma}
\begin{center}
\begin{tabularx}{\textwidth}{L L l}
\toprule
\textbf{Indicador} & \textbf{Fórmula} & \textbf{Leitura} \\
\midrule
Índice de Desempenho de Prazo & Duração planejada concluída até a data de corte $\div$ Duração planejada total até essa data $\times$ 100 & Maior é melhor \\
Tendência da folga do caminho crítico & Folga total medida a cada check-in, comparada ao check-in anterior & Estável ou crescente é melhor; caindo antes do previsto é o alarme antecipado \\
Aderência ao caminho crítico & Atividades críticas concluídas na data planejada $\div$ Total de atividades críticas $\times$ 100 & Maior é melhor \\
Consumo do buffer do projeto & Dias do buffer já utilizados $\div$ Dias totais do buffer $\times$ 100 & Menor, na mesma proporção do avanço do projeto, é melhor \\
\bottomrule
\end{tabularx}
\end{center}

\noindent Esses indicadores são acompanhados no mesmo check-in semanal, com foco no caminho crítico primeiro --- é lá que qualquer atraso vira atraso do projeto inteiro.

\vspace{2mm}
\hrule
\vspace{4mm}

%======================================================================
\section{Prevenção das Armadilhas do Cronograma}
%======================================================================

\noindent As pegadinhas da aula não estão nos números, estão no comportamento da dupla diante do prazo. Cada uma tem uma decisão concreta no TrackFleet:

\begin{itemize}[itemsep=5pt]
  \item \textbf{Lei de Parkinson} (o trabalho se expande até preencher o tempo disponível): ciclos de entrega de 1 semana, com status binário --- pronto (integrado e testado) ou não pronto, sem ``99\% pronto''. A folga não fica dentro de cada atividade; vai para o buffer único do projeto (Seção 5.2).
  \item \textbf{Síndrome do Estudante} (começar tudo na última hora): o check-in semanal de 10 minutos (Governança) cobra status de toda atividade em andamento, não só das que vencem naquela semana --- o desvio aparece na semana 2 do atraso, não na véspera da entrega.
  \item \textbf{Colchões escondidos} (padding/sandbagging): as estimativas das atividades críticas são dadas em três pontos (Seção 3.1), nunca em ponto único ``por segurança''; a reserva de risco fica agregada e visível no buffer do projeto (Seção 5.2), não escondida em cada A6, A7 ou A8.
  \item \textbf{Falácia do Planejamento} (otimismo sistemático ao estimar): as estimativas ancoram na ``visão de fora'' --- comparadas ao tempo real que a dupla levou no Termo de Abertura e no Documento de Escopo (técnica análoga, Seção 3) --- em vez de assumir que ``desta vez vai dar tudo certo''.
  \item \textbf{Síndrome dos 90\%} (os 10\% finais consomem os outros 90\% do tempo): a Definição de Pronto do Documento de Escopo (Seção 3.3) já exige integração e teste com dados reais do rastreador antes de qualquer atividade ser considerada concluída --- ``compila local'' não conta como pronto.
  \item \textbf{Multitarefa e recursos superlotados}: cada atividade tem um único responsável nomeado (Seção 1); o par crítico A6/A7 foi deliberadamente dividido entre Enzo e Kerlon (Seção 4.2) para que nenhum dos dois seja o gargalo do caminho crítico.
\end{itemize}

\end{document}


\chapter{Finanças}\label{cap:financas}
\noindent Aula de 03/09, revisado para o Trabalho Prático 01. Do manual de regras do dinheiro ao valor agregado e ao business case. A planilha \texttt{financas/trackfleet\_financas.xlsx} reproduz todos os cálculos.

\documentclass[11pt,a4paper]{article}
\usepackage[utf8]{inputenc}
\usepackage[T1]{fontenc}
\usepackage[brazil]{babel}
\usepackage[margin=2.5cm]{geometry}
\usepackage{mathptmx}
\usepackage{amsmath}
\usepackage{amssymb}
\usepackage{enumitem}
\usepackage{booktabs}
\usepackage{tabularx}
\usepackage{array}
\usepackage{hyperref}
\hypersetup{colorlinks=false, pdfborder={0 0 0}}
\usepackage{fancyhdr}

\pagestyle{fancy}
\fancyhf{}
\fancyhead[L]{\small TrackFleet --- Finanças}
\fancyhead[R]{\small\thepage}
\renewcommand{\headrulewidth}{0.4pt}

\newcolumntype{L}{>{\raggedright\arraybackslash}X}
\newcolumntype{C}{>{\centering\arraybackslash}X}
\newcolumntype{R}{>{\raggedleft\arraybackslash}X}
\setlist{topsep=3pt}

% Cabeçalho padrão do projeto. No documento consolidado (trabalho_pratico_01/)
% estes três comandos são redefinidos e este mesmo arquivo vira parte do capítulo do domínio.
\newcommand{\cabecalho}[3]{%
  \begin{center}
  {\Large\bfseries DOCUMENTO DE #1 DO PROJETO}\\[2pt]
  {\large\bfseries TrackFleet --- Gestão de Rastreador de Automóveis}\\[4pt]
  Disciplina de Gerenciamento de Projetos $\cdot$ Framework PMBOK\textregistered\ 8\textsuperscript{a} Edição $\cdot$ Domínio: #2\\[2pt]
  Integrantes: Enzo Allebrand e Kerlon Ribeiro (dois GPs) $\cdot$ 4\textsuperscript{o} semestre $\cdot$ Data: #3
  \end{center}
  \vspace{2mm}\hrule\vspace{4mm}}
\newcommand{\parte}[1]{\noindent\textbf{\large #1}\par\vspace{1mm}}
\newcommand{\separador}{\par\vspace{2mm}\hrule\vspace{4mm}}

\begin{document}

\cabecalho{FINANÇAS}{Finanças}{03/09}

\noindent No PMBOK\textregistered\ 8, o domínio de desempenho Finanças cobre quatro processos: Planejar o Gerenciamento Financeiro (o manual de regras do dinheiro do projeto), Estimar Custos (quanto vai custar cada pacote de trabalho), Desenvolver o Orçamento (agregar as estimativas e as reservas em uma linha de base autorizada) e Monitorar e Controlar as Finanças (comparar o realizado com a linha de base e prever onde o projeto vai terminar). A virada de mentalidade da aula guia todo este documento: ficar dentro do orçamento é um resultado de controle; entregar o valor prometido é um resultado estratégico. Por isso o documento não para na planilha de custos --- ele fecha justificando, com ROI, VPL, TIR e payback, por que o dinheiro gasto no TrackFleet produz o resultado prometido.

\noindent Os números desta versão foram recalculados a partir do cronograma nivelado (51 dias úteis) e da reserva de contingência dimensionada pelo registro de riscos. A planilha \texttt{trackfleet\_financas.xlsx} reproduz todos os cálculos com fórmulas.

\separador

%======================================================================
\section{Planejar o Gerenciamento Financeiro: o manual de regras}\label{fin:planejar}
%======================================================================

\noindent Antes de estimar qualquer valor, este processo define como receitas e despesas do TrackFleet serão estimadas, orçadas, monitoradas e controladas. Sem ele, o resto é improviso. As entradas deste plano são o Termo de Abertura, o Documento de Escopo, o Documento de Cronograma e o Documento de Governança; a saída é o conjunto de regras abaixo.

\subsection{Bases de referência}
\begin{itemize}[itemsep=3pt]
  \item \textbf{Moeda:} real brasileiro (R\$); todos os valores deste documento estão em R\$ correntes.
  \item \textbf{Custo-hora de referência:} R\$ 50,00 por hora de trabalho de GP, usado como custo de oportunidade das horas dos dois integrantes --- o projeto não paga salário, mas contabiliza o tempo pelo valor de mercado de uma hora de desenvolvimento júnior, para que o custo real do esforço fique visível e não pareça ``de graça''.
  \item \textbf{Dia útil de trabalho:} 6 horas produtivas; logo, 1 dia-pessoa do cronograma equivale a R\$ 300,00. Essa é a ponte entre o Documento de Cronograma (durações em dias) e este documento (custos em reais).
  \item \textbf{Cadência de medição:} as finanças são medidas no mesmo check-in semanal já definido na Governança e no Cronograma --- não há reunião financeira à parte.
\end{itemize}

\subsection{Fontes de financiamento}
Cada fonte traz regras e nível de rigor diferentes. O TrackFleet combina duas das fontes vistas na aula e descarta as demais:
\begin{itemize}[itemsep=2pt]
  \item \textbf{Orçamento interno (fonte principal):} as horas de trabalho dos dois GPs, no contexto acadêmico da disciplina. É recurso próprio da dupla, sem desembolso de caixa.
  \item \textbf{Contrato/patrocínio de cliente (fonte secundária):} a transportadora parceira (patrocinador) custeia a infraestrutura do piloto --- assinatura da API do provedor de rastreamento, hospedagem e domínio.
  \item \textbf{Subvenções/grants e crowdfunding:} não se aplicam a este projeto.
\end{itemize}
\noindent Como ensina a aula, a fonte define o nível de rigor do controle: por ser financiamento interno e um patrocínio leve de parceiro, a governança financeira é leve (coerente com o modelo híbrido da Governança). A única exceção é a LGPD, que impõe controle formal de tratamento de dados independentemente da fonte.

\subsection{Classificação CapEx {\texttimes} OpEx}
Classificar cada custo na categoria certa é decidido aqui, no planejamento --- classificar errado, como alerta a aula, trava o projeto porque os caminhos de aprovação são diferentes.

\begin{center}
\begin{tabularx}{\textwidth}{l L L}
\toprule
\textbf{Categoria} & \textbf{O que é no TrackFleet} & \textbf{Itens} \\
\midrule
CapEx --- constrói o ativo & Trabalho que cria o ativo de software (a plataforma), gasto de uma vez para gerar um bem duradouro & Horas de trabalho dos GPs em todos os pacotes da EAP, incluindo arquitetura e configuração inicial da infraestrutura \\
OpEx --- roda o dia a dia & Custos contínuos de manter a plataforma no ar durante o projeto e o piloto & Assinatura mensal da API do provedor, hospedagem no VPS, domínio anual, ferramentas; certificado TLS e envio de e-mail em planos gratuitos \\
\bottomrule
\end{tabularx}
\end{center}

\noindent Um ponto concreto liga esta classificação à restrição do Termo de Abertura: o TrackFleet tem orçamento zero para hardware próprio. Na prática, o projeto elimina de propósito o maior CapEx de um sistema de rastreamento --- comprar e instalar rastreadores --- integrando-se por API a dispositivos de terceiros. Esse custo vira OpEx (assinatura da API), muito menor e sem imobilização de capital.

\subsection{Governança financeira e tolerâncias}\label{fin:tolerancias}
Quem aprova o quê segue a matriz RACI da Governança, agora aplicada ao dinheiro:
\begin{itemize}[itemsep=3pt]
  \item \textbf{Gasto do dia a dia (dentro da linha de base):} decidido pelos dois GPs por consenso, sem subir ao patrocinador.
  \item \textbf{Reserva de contingência:} está no controle dos GPs (é para riscos já conhecidos); cada uso é ligado ao risco do registro que o justificou e anotado no log de decisões.
  \item \textbf{Reserva de gerenciamento:} só é liberada pelo patrocinador (Diretoria da transportadora parceira) --- coerente com a aula, que deixa essa reserva com a alta gerência, e com o escalonamento já definido na Governança.
\end{itemize}

\noindent As tolerâncias abaixo são os gatilhos que transformam um número em uma ação --- porque, como diz a aula, relatório colorido parado não é controle:
\begin{center}
\begin{tabularx}{\textwidth}{L L}
\toprule
\textbf{Gatilho} & \textbf{Ação} \\
\midrule
CPI ou SPI abaixo de 0,90 em um check-in & Ação corretiva imediata dos GPs, focada primeiro no caminho crítico \\
Projeção de término (EAC) acima da linha de base, ameaçando a reserva de gerenciamento & Revisão formal e alerta ao patrocinador; acima da reserva, escalonamento e re-baseline \\
Consumo da reserva de contingência acima do avanço físico do projeto & Reavaliação das estimativas restantes e do registro de riscos no check-in seguinte \\
\bottomrule
\end{tabularx}
\end{center}

\separador

%======================================================================
\section{Estimar Custos: quanto vai custar cada pacote}\label{fin:estimar}
%======================================================================

\noindent Estimar custos é uma aproximação periódica do custo dos recursos necessários ao trabalho. A base é a EAP do Documento de Escopo e as atividades do Documento de Cronograma: o custo de cada pacote é a soma dos dias-pessoa das suas atividades, multiplicada por R\$ 300,00. Como na aula, a técnica usada em cada pacote segue a natureza do trabalho.

\subsection{Técnica por pacote}
\begin{itemize}[itemsep=3pt]
  \item \textbf{Análoga:} a gestão do projeto (1.1), que não é uma atividade do cronograma e sim um esforço contínuo (check-ins, relatórios e documentos de cada domínio), é estimada por comparação com o tempo que a dupla já gastou nos documentos das aulas; requisitos (1.2) e encerramento (1.5) seguem as estimativas análogas do Cronograma.
  \item \textbf{Paramétrica:} o Frontend (1.3.3) usa as estimativas paramétricas do Cronograma --- número de telas e de relatórios do catálogo do Documento de Escopo multiplicado pelo tempo médio por unidade.
  \item \textbf{Especialista e bottom-up:} pacotes pequenos e bem conhecidos (autenticação e LGPD, infraestrutura) são estimados somando suas atividades diretamente.
  \item \textbf{Três pontos (PERT):} os três pacotes que contêm as atividades estimadas em três pontos no Cronograma, detalhados na Seção~\ref{fin:pert}.
\end{itemize}

\subsection{Estimativa bottom-up por pacote}\label{fin:pacotes}
Cada pacote é convertido em dias-pessoa a partir das atividades do Cronograma e, pelo custo de R\$ 300,00 por dia-pessoa, em reais. No piloto (A18), cada GP dedica um terço da jornada; A19 e A20 usam os outros dois terços dos seus dias.

\begin{center}
\small
\begin{tabularx}{\textwidth}{l L c c r}
\toprule
\textbf{Pacote EAP} & \textbf{Atividades do Cronograma} & \textbf{Técnica} & \textbf{Dias-pessoa} & \textbf{Custo (R\$)} \\
\midrule
1.1 Gestão do Projeto & Esforço contínuo de gestão (fora da rede) & Análoga & 12,00 & 3.600 \\
1.2 Requisitos e Arquitetura & A1 (dois GPs), A2, A3, A4 & Análoga & 13,00 & 3.900 \\
1.3.1 Backend, API e ingestão & A5, A6, A7 & Três pontos & 15,00 & 4.500 \\
1.3.2 Motor de alertas & A8 & Três pontos & 9,00 & 2.700 \\
1.3.3 Frontend & A11, A12, A13, A14 & Paramétrica & 17,00 & 5.100 \\
1.3.4 Autenticação e LGPD & A9, A10 & Bottom-up & 5,00 & 1.500 \\
1.3.5 Infraestrutura e implantação & A21 & Especialista & 2,00 & 600 \\
1.4 Testes e Homologação & A15, A16, A17, A18 (1/3 de cada GP) & Três pontos & 17,33 & 5.200 \\
1.5 Treinamento e Encerramento & A19 (dois GPs, 2/3), A20 (2/3) & Análoga & 3,33 & 1.000 \\
\midrule
\textbf{Subtotal mão de obra (CapEx)} & & & \textbf{93,67} & \textbf{28.100} \\
Infraestrutura do período (OpEx) & API, VPS, domínio, ferramentas & Bottom-up & --- & 700 \\
\midrule
\textbf{Total das estimativas} & & & & \textbf{28.800} \\
\bottomrule
\end{tabularx}
\end{center}

\subsection{Estimativa de três pontos dos pacotes de risco}\label{fin:pert}
Os três pacotes mais incertos usam os cenários otimista (O), mais provável (M) e pessimista (P) das atividades do Cronograma, convertidos em reais pela mesma fórmula da distribuição Beta:
\[
C_E = \frac{O + 4M + P}{6}
\]

\begin{center}
\begin{tabularx}{\textwidth}{L r r r r r}
\toprule
\textbf{Pacote} & \textbf{O} & \textbf{M} & \textbf{P} & $C_E$ & \textbf{Na linha de base} \\
\midrule
1.3.1 Backend, API e ingestão & 2.700 & 3.900 & 6.600 & 4.150 & 4.500 \\
1.3.2 Motor de alertas & 1.500 & 2.400 & 4.500 & 2.600 & 2.700 \\
1.4 Testes e Homologação & 4.400 & 5.000 & 6.000 & 5.067 & 5.200 \\
\bottomrule
\end{tabularx}
\end{center}

\noindent A coluna ``na linha de base'' é maior que $C_E$ porque o Cronograma arredonda cada $t_E$ para cima (dia inteiro de trabalho). A leitura desses números: o Backend (1.3.1) e o Motor de alertas (1.3.2) carregam a maior distância entre otimista e pessimista --- o Backend depende de uma API externa fora do controle da dupla, e o Motor de alertas reúne quatro lógicas de disparo diferentes na mesma frente. Essa incerteza é tratada no registro de riscos, que dimensiona a reserva de contingência.

\separador

%======================================================================
\section{Desenvolver o Orçamento: da estimativa à linha de base}\label{fin:orcamento}
%======================================================================

\noindent Aqui as estimativas individuais são agregadas e as reservas são adicionadas, formando a linha de base de custos autorizada. Este é o ponto onde mais se erra na prova, como diz a aula: orçamento não é a mesma coisa que linha de base, e as duas reservas entram em lugares diferentes.

\subsection{As duas reservas}
\begin{center}
\begin{tabularx}{\textwidth}{L L L}
\toprule
\textbf{Reserva} & \textbf{Cobre} & \textbf{Onde fica} \\
\midrule
Contingência --- riscos conhecidos & Os riscos do registro de riscos, pelo valor monetário esperado (VME) líquido das ameaças e oportunidades & DENTRO da linha de base; já entra no orçamento inicial e é controlada pelos GPs \\
Gerenciamento --- riscos desconhecidos & Trabalho imprevisto que ninguém mapeou & FORA da linha de base; entra no orçamento total por cima dela e só é liberada pelo patrocinador \\
\bottomrule
\end{tabularx}
\end{center}

\noindent O mnemônico da aula vale para o TrackFleet: contingência para o conhecido.

\subsection{Composição do orçamento}
\begin{center}
\begin{tabularx}{\textwidth}{L r}
\toprule
\textbf{Componente} & \textbf{Valor (R\$)} \\
\midrule
Estimativas dos pacotes (mão de obra) & 28.100 \\
Infraestrutura do período (OpEx) & 700 \\
\textbf{Subtotal das estimativas} & \textbf{28.800} \\
(+) Reserva de contingência (VME líquido do registro de riscos) & 3.970 \\
\textbf{Linha de base de custos (BAC)} & \textbf{32.770} \\
(+) Reserva de gerenciamento (5\% da linha de base) & 1.639 \\
\textbf{Orçamento total do projeto} & \textbf{34.409} \\
\bottomrule
\end{tabularx}
\end{center}

\noindent A linha de base de custos (BAC) é R\$ 32.770: é contra ela que o realizado será medido na Seção~\ref{fin:controlar}. O orçamento total (R\$ 34.409) é maior porque inclui a reserva de gerenciamento --- que existe, mas não faz parte da base de comparação. A reserva de contingência (R\$ 3.970, cerca de 14\% das estimativas) não é um percentual chutado: é a soma dos VMEs do registro de riscos. A linha de base é aprovada pelo patrocinador, na mesma lógica das linhas de base de escopo e de cronograma.

\subsection{Orçamento no tempo: a curva S}\label{fin:curvas}
A linha de base é distribuída no tempo sobre as 11 semanas do Cronograma: cada semana recebe o custo dos dias-pessoa trabalhados nela, uma fração igual da gestão do projeto e do OpEx, e a reserva de contingência proporcional ao custo.

\begin{center}
\footnotesize
\setlength{\tabcolsep}{3pt}
\begin{tabularx}{\textwidth}{l *{11}{C}}
\toprule
\textbf{Semana} & \textbf{S1} & \textbf{S2} & \textbf{S3} & \textbf{S4} & \textbf{S5} & \textbf{S6} & \textbf{S7} & \textbf{S8} & \textbf{S9} & \textbf{S10} & \textbf{S11} \\
\midrule
\% acumulado & 11,8 & 20,4 & 32,2 & 44,0 & 55,7 & 67,5 & 77,2 & 84,8 & 89,6 & 93,8 & 100 \\
Acumulado (R\$) & 3.858 & 6.693 & 10.551 & 14.409 & 18.268 & 22.126 & 25.302 & 27.795 & 29.377 & 30.732 & 32.770 \\
\bottomrule
\end{tabularx}
\end{center}

\noindent A curva tem o miolo íngreme e o fim achatado previstos na aula --- no piloto (S9 a S11) a equipe trabalha só um terço da jornada. O arranque, porém, não é lento: com uma equipe fixa de duas pessoas alocadas desde o primeiro dia, o gasto já começa no ritmo de cruzeiro. Essa curva é o Valor Planejado (PV) ao longo do tempo e é a régua do exemplo resolvido da próxima seção.

\separador

%======================================================================
\section{Monitorar e Controlar as Finanças: o valor agregado}\label{fin:controlar}
%======================================================================

\noindent Este processo supervisiona a saúde financeira: acompanha gastos, ajusta previsões e aplica correções. A ferramenta central é o Gerenciamento do Valor Agregado (EVM), começado desde o dia 1 --- porque, como avisa a aula, fazer EVM só aos 60\% do projeto é inútil: a janela de intervenção já fechou.

\subsection{Os quatro valores-base}
\begin{center}
\begin{tabularx}{\textwidth}{l L L}
\toprule
\textbf{Sigla} & \textbf{Nome} & \textbf{O que é no TrackFleet} \\
\midrule
BAC & Orçamento no término & A linha de base de custos: R\$ 32.770 \\
PV & Valor Planejado & Custo orçado do trabalho programado até a data (a curva S da Seção~\ref{fin:curvas}) \\
EV & Valor Agregado & Custo orçado do trabalho realmente concluído: BAC $\times$ \% concluído \\
AC & Custo Real & O que de fato foi consumido para entregar esse trabalho \\
\bottomrule
\end{tabularx}
\end{center}

\subsection{Variações, índices e previsões}\label{fin:formulas}
\begin{center}
\begin{tabularx}{\textwidth}{L L L}
\toprule
\textbf{Indicador} & \textbf{Fórmula} & \textbf{Leitura} \\
\midrule
CV --- Variação de Custo & $EV - AC$ & Positivo: abaixo do orçamento. Negativo: acima \\
SV --- Variação de Prazo & $EV - PV$ & Positivo: adiantado. Negativo: atrasado \\
CPI --- Índice de Custo & $EV \div AC$ & Acima de 1,0 é bom; abaixo é ruim \\
SPI --- Índice de Prazo & $EV \div PV$ & Acima de 1,0 é bom; abaixo é ruim \\
EAC --- Estimativa no Término & $BAC \div CPI$ & Custo final previsto se a tendência atual continuar \\
ETC --- Estimativa para Terminar & $EAC - AC$ & Quanto ainda falta gastar \\
VAC --- Variação no Término & $BAC - EAC$ & Negativo: projeto termina acima da linha de base \\
TCPI --- Índice para Terminar & $(BAC - EV) \div (BAC - AC)$ & Eficiência exigida do trabalho restante; acima de 1,0 é difícil \\
\bottomrule
\end{tabularx}
\end{center}

\noindent Duas armadilhas da aula ficam registradas aqui: termina em ``Variação'' subtrai, termina em ``Índice'' divide; e SV/SPI medem prazo, mas em unidades de dinheiro (o valor em R\$ do trabalho), nunca em dias.

\subsection{Exemplo resolvido: check-in do fim da Semana 6}\label{fin:exemplo}
Cenário de acompanhamento no corte da Semana 6 (09/10), quando a curva S prevê 67,5\% da linha de base. Suponha que, na medição, o trabalho realmente concluído seja 62\% e o custo real acumulado seja R\$ 21.200:

\begin{center}
\begin{tabularx}{\textwidth}{L L L}
\toprule
\textbf{Valor} & \textbf{Cálculo} & \textbf{Resultado} \\
\midrule
PV & curva S no fim da S6 & R\$ 22.126 \\
EV & $32.770 \times 62\%$ & R\$ 20.317 \\
AC & medido no check-in & R\$ 21.200 \\
CV & $20.317 - 21.200$ & $-$R\$ 883 (acima do orçamento) \\
SV & $20.317 - 22.126$ & $-$R\$ 1.809 (atrasado) \\
CPI & $20.317 \div 21.200$ & 0,96 \\
SPI & $20.317 \div 22.126$ & 0,92 \\
EAC & $32.770 \div (20.317 \div 21.200)$ & R\$ 34.194 \\
ETC & $34.194 - 21.200$ & R\$ 12.994 \\
VAC & $32.770 - 34.194$ & $-$R\$ 1.424 \\
TCPI & $(32.770 - 20.317) \div (32.770 - 21.200)$ & 1,08 \\
\bottomrule
\end{tabularx}
\end{center}

\noindent \textbf{Veredito e decisão:} no ritmo atual o TrackFleet está levemente acima do orçamento (CPI 0,96) e atrasado (SPI 0,92), e a projeção de término (EAC R\$ 34.194) fura a linha de base em R\$ 1.424 --- o equivalente a 87\% da reserva de gerenciamento (R\$ 1.639). Pela tolerância da Seção~\ref{fin:tolerancias}, isso ainda não estoura o orçamento total, então não é escalonamento imediato, mas dispara revisão formal e coloca o patrocinador em alerta. A ação corretiva se concentra no maior risco do caminho crítico, o Motor de alertas (atividade A8 do Cronograma). O TCPI de 1,08 diz o preço disso: o trabalho restante precisa ser feito cerca de 8\% mais eficiente do que o planejado --- exigente, mas possível. E, seguindo a dica do professor, a decisão olha a tendência do CPI ao longo dos check-ins, não a foto de uma semana isolada.

\separador

%======================================================================
\section{O Valor: por que o investimento se justifica}\label{fin:valor}
%======================================================================

\noindent Cumprir o orçamento não é sucesso financeiro. Esta seção responde à pergunta executiva da aula --- ``o dinheiro gasto produziu o resultado prometido?'' --- medindo o retorno do TrackFleet para o patrocinador. O valor, como reconhece o padrão, não é só financeiro: é tratado nas duas dimensões.

\subsection{Valor tangível: o business case}
As premissas abaixo são estimativas de business case do ponto de vista da transportadora parceira; são deliberadamente conservadoras e assumidamente frágeis (por isso a análise de sensibilidade adiante). Elas não se confundem com o ``indicador de economia em reais'' que o Documento de Escopo deixou de fora do produto --- aqui o cálculo é do projeto, não uma tela da plataforma.

\begin{itemize}[itemsep=2pt]
  \item \textbf{Investimento inicial:} orçamento total do projeto, R\$ 34.409. É uma visão conservadora: inclui o custo de oportunidade das horas da dupla, que a transportadora não desembolsa.
  \item \textbf{Ganho anual estimado:} R\$ 24.000/ano, pela redução de manutenção corretiva, menos multas por documento vencido e menor uso indevido dos veículos.
  \item \textbf{Custo operacional anual (pós-projeto):} R\$ 3.600/ano de OpEx (API, hospedagem e domínio).
  \item \textbf{Ganho líquido anual:} R\$ 24.000 $-$ R\$ 3.600 $=$ R\$ 20.400/ano.
  \item \textbf{Taxa de desconto (custo de capital):} 12\% ao ano.
\end{itemize}

\begin{center}
\begin{tabularx}{\textwidth}{L L L}
\toprule
\textbf{Métrica} & \textbf{Cálculo} & \textbf{Resultado} \\
\midrule
Payback & $34.409 \div 20.400$ & $\approx$ 1,7 ano ($\approx$ 20 meses) \\
VPL (3 anos) & $-34.409 + \sum_{t=1}^{3} \dfrac{20.400}{1{,}12^{t}}$ & $\approx +$R\$ 14.588 \\
TIR (3 anos) & taxa que zera o VPL & $\approx$ 35\% ao ano \\
ROI (3 anos) & $(3 \times 20.400 - 34.409) \div 34.409$ & $\approx$ 78\% \\
\bottomrule
\end{tabularx}
\end{center}

\noindent Leitura: o VPL é positivo (o projeto cria valor mesmo descontando o dinheiro no tempo), a TIR de cerca de 35\% supera com folga o custo de capital de 12\% (logo, aprova-se o investimento) e o payback fica abaixo de dois anos.

\noindent Sensibilidade: mesmo que o ganho anual seja 30\% menor que o previsto (R\$ 16.800 em vez de R\$ 24.000), o ganho líquido cai para R\$ 13.200/ano e o payback estica para cerca de 2,6 anos. O VPL de 3 anos fica negativo (cerca de $-$R\$ 2.705) e volta a positivo em um horizonte de 4 anos (cerca de $+$R\$ 5.684): a decisão de investir se sustenta, com margem mais estreita e dependendo de a plataforma ser usada por mais tempo.

\subsection{Valor intangível}
Nem todo valor cabe na planilha. O TrackFleet entrega também:
\begin{itemize}[itemsep=2pt]
  \item \textbf{Conformidade regulatória:} adequação à LGPD no tratamento de dados de localização e de motoristas --- valor obrigatório, não negociável (VBS do Documento de Escopo).
  \item \textbf{Segurança:} maior proteção de veículos e motoristas com o monitoramento e os alertas.
  \item \textbf{Decisão baseada em dados:} gestão da frota apoiada em dados reais, não em achismo --- o valor esperado do Termo de Abertura.
  \item \textbf{Satisfação do cliente:} uma solução acessível para um segmento hoje mal atendido pelo custo elevado das plataformas corporativas.
  \item \textbf{Aprendizado da dupla:} domínio prático do PMBOK\textregistered\ 8 aplicado de ponta a ponta, valor acadêmico do projeto.
\end{itemize}

\subsection{Alinhamento com o valor do produto}
O business case reforça a priorização já registrada na Value Breakdown Structure do Documento de Escopo: Monitoramento em tempo real (E1) e Alertas automáticos (E2), que concentram 70\% do valor, são exatamente os entregáveis que produzem a economia de manutenção corretiva e de uso indevido usada no cálculo acima. Se o prazo do semestre apertar, cortar relatórios avançados preserva o retorno; cortar E1 ou E2 destruiria o business case.

\separador

%======================================================================
\section{Armadilhas Financeiras e como o TrackFleet as evita}\label{fin:armadilhas}
%======================================================================

\noindent As pegadinhas da aula não estão nos números, e sim no comportamento diante deles. Cada uma tem uma decisão concreta no projeto:

\begin{itemize}[itemsep=5pt]
  \item \textbf{Confundir linha de base com orçamento:} a reserva de gerenciamento (R\$ 1.639) fica de fora da linha de base (R\$ 32.770) e só entra no orçamento total (R\$ 34.409). A comparação do EVM é sempre contra a linha de base, nunca contra o orçamento total.
  \item \textbf{Trocar índice por variação:} a tabela de fórmulas (Seção~\ref{fin:formulas}) registra que termina em ``Índice'' divide e termina em ``Variação'' subtrai --- e que SV/SPI medem prazo em reais, não em dias.
  \item \textbf{Inverter a leitura do TCPI:} para CPI e SPI, abaixo de 1,0 é ruim; para o TCPI, o oposto --- acima de 1,0 significa que o trabalho restante precisa ser mais eficiente, como no exemplo da Seção~\ref{fin:exemplo} (1,08).
  \item \textbf{Achar que orçamento cumprido é sucesso:} por isso este documento não para no controle de custos --- a Seção~\ref{fin:valor} mede o valor entregue (ROI, VPL, TIR, payback), que é o resultado estratégico.
  \item \textbf{Começar o EVM tarde demais:} a medição começa no dia 1, no mesmo check-in semanal, para que o desvio apareça enquanto a janela de intervenção ainda está aberta.
  \item \textbf{Relatório sem narrativa:} cada indicador do exemplo resolvido vem com a decisão que ele dispara (revisão, ação corretiva, alerta ao patrocinador) --- dado sem decisão não é controle.
  \item \textbf{Escolher o EAC errado:} como o desvio do TrackFleet é uma variação típica (tendência de custo), o EAC usa $BAC \div CPI$. A fórmula de evento único, $AC + (BAC - EV)$, só valeria para um desvio isolado que não se repete; e a fórmula pessimista, $AC + (BAC - EV) \div (CPI \times SPI)$, fica para quando custo e prazo pesam juntos sobre o restante do projeto.
\end{itemize}

\end{document}


\chapter{Partes Interessadas}\label{cap:partes}
\noindent Aula de 10/09. Quem são as partes interessadas, como cada uma é engajada e como a informação flui entre elas.
\subdocumento{../partes_interessadas/registro_das_partes_interessadas}
\subdocumento{../partes_interessadas/plano_de_engajamento}
\subdocumento{../partes_interessadas/plano_de_comunicacoes}

\chapter{Recursos}\label{cap:recursos}
\noindent Aula de 17/09. A equipe e os recursos físicos e virtuais: acordos da dupla, responsabilidades, estrutura dos recursos e carga semanal.
\subdocumento{../recursos/termo_da_equipe}
\subdocumento{../recursos/matriz_raci}
\subdocumento{../recursos/estrutura_analitica_dos_recursos}
\subdocumento{../recursos/estimativa_e_histograma}

\chapter{Riscos}\label{cap:riscos}
\noindent Aula de 24/09. Como os riscos são gerenciados, quais são, quanto valem e como são respondidos e monitorados.
\subdocumento{../riscos/plano_de_gerenciamento_dos_riscos}
\subdocumento{../riscos/registro_de_riscos}
\subdocumento{../riscos/analise_e_respostas}

%======================================================================
% Integração
%======================================================================
\capitulosemnumero{Integração entre os Domínios}

\noindent Governança, como mostrou a primeira aula, não é um portão: é a espinha dorsal que conecta os outros seis domínios. No TrackFleet essa conexão é feita por números e decisões que aparecem em mais de um capítulo e precisam bater entre si.

\secaosemnumero{Como os Domínios se Amarram}
\begin{center}
\small
\begin{tabularx}{\textwidth}{>{\raggedright\arraybackslash}p{3.6cm} L}
\toprule
\textbf{Ligação} & \textbf{Como aparece no projeto} \\
\midrule
Governança $\to$ Cronograma e Recursos & A autoridade técnica dividida por frente na RACI de Governança é a mesma da RACI de Recursos e define os donos das atividades do Cronograma. \\
Escopo $\to$ Cronograma & Cada pacote da EAP vira atividades (A1 a A21); a Definição de Pronto vira o status binário cobrado no check-in. \\
Cronograma $\to$ Finanças & Os dias-pessoa das atividades, a R\$ 300 cada, formam os custos; as 11 semanas do calendário formam a curva S do valor planejado. \\
Riscos $\to$ Finanças & O VME líquido do registro de riscos é a reserva de contingência (R\$ 3.970) dentro da linha de base. \\
Riscos $\to$ Cronograma & O buffer de 24 dias úteis protege o prazo; o limite de 5 dias de atraso no caminho crítico aciona o patrocinador. \\
Recursos $\to$ Cronograma & O histograma prova que o cronograma nivelado não sobrecarrega ninguém; a primeira linha de base sobrecarregava. \\
Partes Interessadas $\to$ Riscos & O patrocinador (R06), os motoristas (R07) e o provedor da API (R01) são partes interessadas e também fontes de risco, com a mesma pessoa como dono no engajamento e no risco. \\
Finanças $\to$ Escopo & O business case depende dos épicos E1 e E2, os mesmos que a VBS do Escopo prioriza com 70\% do valor. \\
\bottomrule
\end{tabularx}
\end{center}

\secaosemnumero{O Check-in Semanal como Ponto de Integração}
Os sete domínios são acompanhados no mesmo check-in de 10 minutos, na mesma ordem:
\begin{enumerate}[itemsep=2pt]
  \item \textbf{Escopo e Cronograma:} status binário das atividades em andamento, tendência da folga e consumo do buffer.
  \item \textbf{Finanças:} CPI e SPI frente às tolerâncias.
  \item \textbf{Riscos:} exposição frente à meta, gatilhos disparados e riscos novos.
  \item \textbf{Recursos:} horas planejadas e realizadas de cada GP.
  \item \textbf{Partes Interessadas:} nível de engajamento nas partes com lacuna.
  \item \textbf{Governança:} decisões tomadas, registradas no log, e o que precisa subir ao patrocinador.
\end{enumerate}

\secaosemnumero{Uso de Inteligência Artificial}
Como pedem as aulas de Partes Interessadas e de Riscos, a dupla documentou o uso de IA nas seções correspondentes (Capítulos~\ref{cap:partes} e~\ref{cap:riscos}). Em resumo: um assistente de IA apoiou a revisão de consistência entre os documentos, a conferência dos cálculos e a redação de rascunhos; as decisões do projeto e a validação de cada número ficaram com os GPs, e nenhum dado pessoal real foi enviado à ferramenta.

\secaosemnumero{Próximos Passos}
O projeto segue em execução com o check-in semanal. As próximas aulas da disciplina tratam de métodos ágeis (Scrum, Kanban, XP), que se aplicam diretamente ao backlog do produto do Capítulo~\ref{cap:escopo}: os épicos e histórias já definidos são a entrada natural para organizar o trabalho em sprints e em um quadro Kanban.

\end{document}
